
\subsubsection{The Training-Generation Gap}

The central finding is negative: BiDPO training creates gradient pressure toward collaboration patterns during optimization, but this pressure does not transfer to generation-time behavior at proof-of-concept scale.

Several factors may contribute to this gap:

\begin{enumerate}
\item \textbf{Insufficient training duration.} The 250-step proof-of-concept scale may be insufficient to induce behavioral change. Longer training (full epoch, ~10,000 steps) may show different results.
\end{enumerate}

\begin{enumerate}
\item \textbf{Suboptimal $\lambda$ weighting.} Only $\lambda$ = 0.5 was tested. The agency signal may require stronger weighting ($\lambda$ > 0.5) to compete with the preference signal, or weaker weighting to avoid interference.
\end{enumerate}

\begin{enumerate}
\item \textbf{Heuristic limitations.} The pattern-matching collaboration score may capture surface features that do not correspond to behaviors the model can learn to produce at generation time.
\end{enumerate}

\begin{enumerate}
\item \textbf{Base model dominance.} Mistral-7B-Instruct has strong priors from instruction tuning that 250 steps on 4000 samples cannot override.
\end{enumerate}

\subsubsection{Limitations}

\begin{itemize}
\item \textbf{Proof-of-concept scale only:} All training used 250 steps on 4000 samples. Full-scale training (~10,000 steps on 170,000 samples) was not conducted.
\end{itemize}

\begin{itemize}
\item \textbf{Single model architecture:} Only Mistral-7B-Instruct-v0.2 was tested. Results may not generalize to other model sizes or architectures.
\end{itemize}

\begin{itemize}
\item \textbf{Single $\lambda$ value:} Only $\lambda$ = 0.5 was tested. The optimal weighting remains unknown.
\end{itemize}

\begin{itemize}
\item \textbf{Heuristic-based scoring:} The collaboration score uses regex pattern matching rather than learned representations. Heuristics may not capture the full range of agency-preserving behaviors.
\end{itemize}

\begin{itemize}
\item \textbf{Downstream benchmarks not evaluated:} MT-Bench and TruthfulQA evaluations were planned but not executed due to the E3 failure.
\end{itemize}

\subsubsection{Value of Negative Results}

This negative result contributes by: (1) demonstrating that the foundational mechanism is feasible (orthogonality and stability), (2) identifying the specific failure point (generation transfer at proof-of-concept scale), and (3) informing future work with concrete next steps rather than open-ended exploration.

